\section{Porovnání možností nasazení}
\label{sec:selection-deployment}

Jazykové modely lze do systému integrovat dvěma základními způsoby: prostřednictvím cloudového API nebo lokálním on-premise nasazením.
Každý z přístupů s sebou nese vlastní sadu výhod, omezení a kompromisů.
Tato sekce analyzuje obě varianty z hlediska hardwarových nároků, provozních limitů, ochrany dat a nákladů.

\subsection{Cloudové API}
\label{subsec:deployment-cloud}

Cloudové API je způsob přístupu k jazykovému modelu provozovanému na serverech externího poskytovatele.
Integrace probíhá prostřednictvím standardních HTTP volání -- vývojář odešle dotaz (prompt) na API endpoint poskytovatele a obdrží odpověď modelu.
Výpočet probíhá plně na straně poskytovatele; provozovatel aplikace (v tomto případě systém Kelvin) nepotřebuje žádný specializovaný hardware.

Mezi nejznámější poskytovatele cloudových API pro jazykové modely patří:
\begin{itemize}
    \item \textbf{OpenAI} -- poskytuje přístup k modelům řady GPT-4o a GPT-4o-mini přes API \cite{openai_api}.
    \item \textbf{Anthropic} -- poskytuje přístup k modelům Claude~3.5 Sonnet a Claude~3.5 Haiku \cite{anthropic_api}.
    \item \textbf{Google} -- přístup k modelům Gemini prostřednictvím Google AI Studio nebo Vertex AI \cite{google_vertex}.
    \item \textbf{Mistral AI} -- poskytuje přístup k vlastním modelům Mistral Large a Mistral Nemo, které jsou k dispozici i jako open-source\footnote{Open-source zde znamená, že váhy modelu jsou volně dostupné ke stažení. Open-source LLM lze nasadit lokálně bez nutnosti platit za API přístupy.} ke stažení \cite{mistral2023}.
\end{itemize}

\subsubsection*{Hardwarové nároky}

Z pohledu provozovatele systému Kelvin jsou hardwarové nároky cloudového API prakticky nulové.
Veškerý výpočet probíhá na serverech poskytovatele; na straně Kelvinu je potřeba pouze stabilní internetové připojení a server, který zvládne odesílat HTTP požadavky na API endpoint.
To je výrazná výhoda oproti on-premise nasazení, které vyžaduje specializovaný hardware.

\subsubsection*{Limity}

Cloudová API zavádějí několik typů omezení, která je nutné při návrhu systému zohlednit:

\begin{itemize}
    \item \textbf{Limit počtu tokenů v kontextu} -- každý model má maximální délku zpracovatelného vstupu i výstupu.
    Toto omezení je zpravidla dostatečné (běžně 128\,000 tokenů nebo více), ale je třeba jej sledovat pro případ velmi rozsáhlých odevzdání.

    \item \textbf{Rate limiting}\footnote{Rate limiting (omezení frekvence požadavků) je mechanismus, jímž poskytovatel API omezuje, kolik požadavků může klient odeslat za danou časovou jednotku. Typicky se udává jako RPM -- requests per minute.}
    -- poskytovatelé omezují počet požadavků za minutu (RPM) a počet tokenů za minutu (TPM) pro každý API klíč.
    V případě hromadného odevzdávání úloh před deadlinem, kdy mnoho studentů odešle řešení v krátkém časovém okně, může systém narazit na tyto limity a být nucen požadavky zpomalit nebo zařadit do fronty.

    \item \textbf{Závislost na dostupnosti poskytovatele} -- pokud API poskytovatele není dostupné (výpadek, plánovaná odstávka), systém nemůže provádět analýzu.
    Tento bod je zmírněn skutečností, že analýza v systému Kelvin probíhá asynchronně a není nutně vyžadována okamžitě při odevzdání.
\end{itemize}

\subsubsection*{Ochrana dat}

Použití cloudového API znamená, že vstupní data -- mimo jiné zdrojové kódy studentů a zadání úloh -- jsou odesílána na servery externího poskytovatele a zde zpracovávána.
To vyvolává otázky z hlediska ochrany osobních údajů a souladu s nařízením GDPR \cite{gdpr2016}.

Klíčové otázky, které je nutné s poskytovatelem API vyřešit:
\begin{itemize}
    \item Jsou vstupní data uchovávána po zpracování, a pokud ano, jak dlouho?
    \item Jsou data využívána pro trénování modelů? (Toto je pro ochranu soukromí studentů obzvláště citlivé.)
    \item Kde fyzicky probíhá zpracování dat? Je zpracování realizováno na serverech v EHP?
    \item Je poskytovatel ochoten uzavřít smlouvu o zpracování osobních údajů (Data Processing Agreement, DPA)?
\end{itemize}

Velcí poskytovatelé jako OpenAI a Anthropic nabízejí zákazníkům uzavření DPA a mají varianty API, u nichž se zavazují, že data nebudou použita pro trénování modelů.
Záleží však na konkrétním smluvním uspořádání.
Pro českou veřejnou vysokou školu tak cloudové API není automaticky nevhodné, ale vyžaduje právní ošetření ze strany institutu nebo školy \cite{gdpr2016}.

\subsubsection*{Náklady}

Náklady cloudových API jsou určeny cenovým modelem \emph{za token}.
Poskytovatelé zpravidla rozlišují vstupní tokeny (text odesílaný modelu) a výstupní tokeny (text generovaný modelem), přičemž výstupní tokeny jsou typicky dražší.

\begin{table}[H]
    \centering
    \resizebox{\textwidth}{!}{%
        \begin{tabular}{llll}
            \toprule
            \textbf{Model}              & \textbf{Vstupní tokeny (\$/1M)} & \textbf{Výstupní tokeny (\$/1M)} & \textbf{Přibližný náklad / odevzdání} \\
            \midrule
            GPT-4o (OpenAI)             & 2,50                            & 10,00                            & $\sim$0,01--0,05 \$ \\
            GPT-4o-mini (OpenAI)        & 0,15                            & 0,60                             & $\sim$0,001--0,005 \$ \\
            Claude 3.5 Sonnet           & 3,00                            & 15,00                            & $\sim$0,01--0,06 \$ \\
            Claude 3.5 Haiku            & 0,80                            & 4,00                             & $\sim$0,003--0,015 \$ \\
            Gemini 1.5 Pro              & 1,25                            & 5,00                             & $\sim$0,005--0,025 \$ \\
            Gemini 1.5 Flash            & 0,075                           & 0,30                             & $\sim$0,0003--0,001 \$ \\
            \bottomrule
        \end{tabular}
    }
    \caption{Přibližné ceny vybraných cloudových API modelů (stav 2024--2025) \cite{openai_api, anthropic_api, google_vertex}. Náklady na jedno odevzdání jsou orientační výpočet pro vstup 4\,000 tokenů a výstup 1\,000 tokenů.}
    \label{tab:cloud-api-costs}
\end{table}

Pro systém se stovkami odevzdání ročně jsou náklady za použití cenově dostupnějších modelů (GPT-4o-mini, Claude Haiku, Gemini Flash) velmi přijatelné -- v řádu jednotek dolarů za semestr.
Při použití nejvýkonnějších modelů (GPT-4o, Claude Sonnet) mohou náklady vzrůst, avšak stále zůstávají v rozumných mezích pro akademické použití.

\subsection{On-premise řešení}
\label{subsec:deployment-on-premise}

On-premise nasazení znamená provozování jazykového modelu přímo na serverech instituce, bez využití externích cloudových služeb.
Pro tuto variantu existuje několik softwarových nástrojů usnadňujících správu a inferenci modelů:

\begin{itemize}
    \item \textbf{Ollama} \cite{ollama} -- nástroj umožňující snadné stažení a spuštění open-source jazykových modelů lokálně.
    Poskytuje REST API kompatibilní s rozhraním OpenAI, což zjednodušuje migraci mezi cloudovým a lokálním nasazením.

    \item \textbf{vLLM} \cite{vllm2023} -- výkonný inference server pro nasazení v produkčním prostředí.
    Implementuje optimalizace jako \emph{PagedAttention}\footnote{PagedAttention je technika správy paměti pro KV cache, která ji rozděluje do stránek pevné velikosti podobně jako virtuální paměť v operačních systémech. To umožňuje efektivnější využití paměti GPU a výrazně vyšší propustnost při obsluze více souběžných požadavků.}, díky nimž dosahuje výrazně vyšší propustnosti než naivní implementace.
    Je vhodný pro produkční nasazení s více souběžnými uživateli.

    \item \textbf{llama.cpp} \cite{llamacpp} -- implementace inference v jazyce C++, schopná provozovat modely i na CPU nebo na méně výkonných GPU.
    Podporuje kvantizaci\footnote{Kvantizace je technika komprese modelu, při níž jsou váhy sítě reprezentovány s menší přesností (například 4 nebo 8 bitů namísto 16). Tím se výrazně sníží paměťové nároky při zachování přijatelné kvality výstupů.} do formátů GGUF, díky níž lze modely provozovat na hardware s omezenou pamětí GPU.
\end{itemize}

\subsubsection*{Hardwarové nároky}

Pro on-premise nasazení je zcela zásadní dostupnost vhodného hardware.
Inference jazykových modelů je výpočetně náročná operace (viz sekce~\ref{subsec:llm-computational-cost}), proto je využití GPU\footnote{GPU (Graphics Processing Unit) -- grafický procesor s tisíci paralelními výpočetními jádry, vhodný pro maticové operace tvořící základ LLM inference.}
s dostatečnou kapacitou paměti VRAM\footnote{VRAM (Video RAM) -- specializovaná paměť grafického procesoru, do níž musí být uloženy váhy celého modelu. Nedostatečná VRAM znemožňuje spuštění modelu.}
téměř nezbytné pro přijatelné výkony.

\begin{table}[H]
    \centering
    \resizebox{\textwidth}{!}{%
        \begin{tabular}{lllll}
            \toprule
            \textbf{Model (velikost)} & \textbf{Přesnost FP16} & \textbf{Kvantizace Q8} & \textbf{Kvantizace Q4} & \textbf{Minimální GPU} \\
            \midrule
            7B parametrů   & $\sim$14 GB & $\sim$7 GB  & $\sim$4--5 GB & RTX 3080 / 4070 (10--12 GB) \\
            13B parametrů  & $\sim$26 GB & $\sim$13 GB & $\sim$7--9 GB & RTX 3090 / 4090 (24 GB) \\
            32B parametrů  & $\sim$64 GB & $\sim$32 GB & $\sim$18 GB   & A100 (40 GB) nebo 2× RTX 3090 \\
            70B parametrů  & $\sim$140 GB & $\sim$70 GB & $\sim$40 GB  & A100 80 GB nebo více GPU \\
            \bottomrule
        \end{tabular}
    }
    \caption{Přibližné hardwarové nároky open-source modelů pro on-premise nasazení. Hodnoty jsou orientační a závisí na konkrétní architektuře modelu. Nezahrnují paměť nutnou pro KV cache.}
    \label{tab:on-premise-hardware}
\end{table}

\subsubsection*{Limity}

On-premise nasazení přináší vlastní sadu omezení:

\begin{itemize}
    \item \textbf{Omezení dostupným hardware} -- kvalita a rychlost analýzy jsou přímou funkcí dostupného GPU.
    Na nedostatečném hardware je nutné volit menší nebo silněji kvantizované modely, což může snižovat kvalitu výstupů.

    \item \textbf{Škálovatelnost} -- cloudové API automaticky škáluje s nárůstem zátěže; lokální server má fixní kapacitu.
    Při hromadném odevzdávání se požadavky řadí do fronty a čekají na zpracování.

    \item \textbf{Správa modelu} -- nové verze modelů musí být ručně staženy a nasazeny.
    Aktualizace nebo výměna modelu vyžaduje přímou intervenci správce systému.
\end{itemize}

\subsubsection*{Ochrana dat}

Z pohledu ochrany dat je on-premise nasazení výrazně jednodušší situace.
Studentské zdrojové kódy a zadání úloh nikdy neopustí infrastrukturu školy.
Neexistuje třetí strana, které by musela být svěřena správa dat, a není potřeba uzavírat DPA s externím poskytovatelem.
Z hlediska GDPR jde o přímočarou variantu: správce dat (škola) je totožný s provozovatelem infrastruktury.

\subsubsection*{Náklady}

Náklady on-premise nasazení mají odlišnou strukturu než u cloudového API:

\begin{itemize}
    \item \textbf{Jednorázové pořizovací náklady} -- GPU server vhodný pro inferenci (např. NVIDIA RTX 4090 s 24 GB VRAM) se pohybuje v ceně od přibližně 60\,000 Kč výše.
    Pro výkonnější varianty (NVIDIA A100) jsou náklady řádově vyšší (200\,000--500\,000 Kč).

    \item \textbf{Průběžné provozní náklady} -- elektřina (moderní GPU má spotřebu 150--400 W při zatížení), chlazení serverovny a lidská práce nutná pro správu a aktualizaci systému.

    \item \textbf{Nulové náklady za token} -- po zaplacení hardware jsou marginalní náklady na každé zpracované odevzdání prakticky nulové, bez ohledu na počet odevzdání.
\end{itemize}

Pro systém s malým počtem uživatelů a odevzdání může být on-premise pořizovací investice ekonomicky nevýhodná ve srovnání s cloudovým API.
Naopak pro systém se stovkami studentů zpracovávajícími tisíce odevzdání ročně se on-premise investice může vyplatit v horizontu 1--2 let.

\begin{table}[H]
    \centering
    \resizebox{\textwidth}{!}{%
        \begin{tabular}{lll}
            \toprule
            \textbf{Vlastnost}           & \textbf{Cloudové API}                     & \textbf{On-premise nasazení} \\
            \midrule
            Hardwarové nároky            & Žádné (výpočet u poskytovatele)           & GPU server (10+ GB VRAM) \\
            Pořizovací náklady           & Nulové                                    & Vysoké (GPU hardware) \\
            Průběžné náklady             & Za token (opakující se)                   & Elektřina a správa \\
            Škálovatelnost               & Automatická (rate limity)                 & Fixní kapacita \\
            Ochrana dat                  & Data opouštějí školu (DPA nutné)          & Data zůstávají na škole \\
            Kvalita modelů               & Nejnovější (GPT-4o, Claude Sonnet)        & Závisí na hardware \\
            Latence                      & Závislá na internetu a zátěži API         & Závisí na hardware \\
            Správa a aktualizace         & Automatická (poskytovatel)                & Manuální (správce) \\
            \bottomrule
        \end{tabular}
    }
    \caption{Srovnání cloudového API a on-premise nasazení jazykového modelu z hlediska klíčových vlastností}
    \label{tab:deployment-comparison}
\end{table}

\endinput
